\documentclass[11pt, a4paper]{article}

% =====================================================================
% PACKAGES
% =====================================================================
\usepackage[utf8]{inputenc}
\usepackage[T1]{fontenc}
\usepackage{lmodern}
\usepackage{geometry}
\usepackage{hyperref}
\usepackage{xcolor}
\usepackage{listings}
\usepackage{amsmath}
\usepackage{graphicx}
\usepackage{titlesec}
\usepackage{parskip}
\usepackage{booktabs}
\usepackage{enumitem}

\geometry{a4paper, left=25mm, right=25mm, top=25mm, bottom=25mm}

% =====================================================================
% COLOR THEME
% =====================================================================
\definecolor{codegreen}{rgb}{0,0.55,0}
\definecolor{codegray}{rgb}{0.4,0.4,0.4}
\definecolor{codepurple}{rgb}{0.5,0,0.75}
\definecolor{codeblue}{rgb}{0.0,0.2,0.75}
\definecolor{backcolour}{rgb}{0.96,0.96,0.94}
\definecolor{accentblue}{rgb}{0.1,0.2,0.6}

% =====================================================================
% LISTINGS STYLES
% =====================================================================
\lstdefinestyle{cppstyle}{
    backgroundcolor=\color{backcolour},
    commentstyle=\color{codegreen}\itshape,
    keywordstyle=\color{codeblue}\bfseries,
    numberstyle=\tiny\color{codegray},
    stringstyle=\color{codepurple},
    basicstyle=\ttfamily\footnotesize,
    breaklines=true, captionpos=b, keepspaces=true,
    numbers=left, numbersep=6pt,
    showspaces=false, showstringspaces=false,
    tabsize=2, frame=single, rulecolor=\color{codegray},
    language=C++
}

\lstdefinestyle{tsstyle}{
    backgroundcolor=\color{backcolour},
    commentstyle=\color{codegreen}\itshape,
    keywordstyle=\color{codeblue}\bfseries,
    numberstyle=\tiny\color{codegray},
    stringstyle=\color{codepurple},
    basicstyle=\ttfamily\footnotesize,
    breaklines=true, captionpos=b, keepspaces=true,
    numbers=left, numbersep=6pt,
    showspaces=false, showstringspaces=false,
    tabsize=2, frame=single, rulecolor=\color{codegray},
    language=Java,
    morekeywords={const,let,interface,export,from,import,type,
                  useState,useEffect,useRef,useDispatch,useSelector,
                  PayloadAction,createSlice,configureStore,React}
}

\lstdefinestyle{bashstyle}{
    backgroundcolor=\color{backcolour},
    basicstyle=\ttfamily\footnotesize,
    breaklines=true, captionpos=b,
    frame=single, rulecolor=\color{codegray},
    numbers=none, showstringspaces=false, language=bash
}

\lstset{style=cppstyle}

% =====================================================================
% HYPERREF
% =====================================================================
\hypersetup{
    colorlinks=true, linkcolor=accentblue, urlcolor=accentblue,
    pdftitle={Raft Consensus Simulator: A Step-by-Step Guide},
}

% =====================================================================
% SECTION FORMATTING
% =====================================================================
\titleformat{\section}[block]
  {\Large\bfseries\color{accentblue}}{\thesection}{1em}{}[\titlerule]
\titleformat{\subsection}[block]
  {\large\bfseries}{\thesubsection}{1em}{}

% =====================================================================
% DOCUMENT
% =====================================================================
\begin{document}

\begin{titlepage}
    \centering
    \vspace*{3cm}
    {\Huge\bfseries\color{accentblue} Raft Consensus Simulator\par}
    \vspace{0.5cm}
    {\Large A Step-by-Step Implementation Guide\par}
    \vspace{1cm}
    \rule{0.6\textwidth}{1pt}
    \vspace{1cm}
    {\large A comprehensive walkthrough of building a full-stack,
    real-time distributed systems simulator from scratch, covering
    C++17 multithreading, WebSocket telemetry, and React state
    visualization.\par}
    \vfill
    {\large \today\par}
\end{titlepage}

\tableofcontents
\newpage

% =====================================================================
\section{Step 1: Project Architecture and Technology Selection}
% =====================================================================

\subsection{Conceptual Overview}
Before writing a single line of code, the most consequential engineering
decisions must be made: \textit{what to build, with what tools, and why}.
The Raft Consensus Simulator is a full-stack application with two
fundamentally distinct concerns. The first is the \textbf{consensus
engine}---a correct, high-performance implementation of the Raft
algorithm---which demands deterministic timing and direct control over
threads. The second is the \textbf{visualization layer}---a real-time UI
that renders the cluster's live state---which demands a reactive,
component-driven architecture.

The central design insight is that these two concerns must be rigorously
decoupled. If the visualization logic were embedded in the consensus engine,
or vice versa, the system would be fragile and impossible to test. The
solution is a clean, bidirectional bridge: a WebSocket connection.

\subsection{The Building Process}
The architecture is divided into two independent, runnable services:
\begin{enumerate}[leftmargin=*]
    \item \textbf{Backend} (\texttt{/backend}): A C++17 process housing
          all Raft logic, threading, and a Crow-powered WebSocket server.
    \item \textbf{Frontend} (\texttt{/frontend}): A Vite/React/TypeScript
          SPA that connects to the backend WebSocket and renders live state.
\end{enumerate}

\subsection{Implementation Details}
\textbf{Why C++17?} The Raft algorithm's correctness depends on timing:
election timeouts must fire within precise, randomized windows.
Garbage-collected runtimes (Java, Go) introduce non-deterministic GC pauses
that can trigger spurious elections. C++ guarantees deterministic execution,
and \texttt{std::thread} / \texttt{std::mutex} provide the exact primitives
needed to map each Raft node onto a dedicated OS thread.

\textbf{Why React + Redux Toolkit?} The frontend receives 30+ state updates
per second during elections. Naively piping these into \texttt{useState}
would trigger a full component-tree re-render on every message. Redux
Toolkit, backed by Immer.js, merges incoming data into an immutable global
state tree; React's diffing algorithm then re-renders only the components
whose data actually changed.

\textbf{Why WebSockets?} HTTP polling introduces latency and wasted
bandwidth. WebSockets establish a persistent, full-duplex TCP tunnel: the
backend \textit{pushes} telemetry events the moment they occur, and the
frontend can \textit{push} fault-injection commands back through the same
connection.

\subsection{Code \& Configuration}
\begin{lstlisting}[style=cppstyle, caption={backend/CMakeLists.txt}]
cmake_minimum_required(VERSION 3.14)
project(RaftSimulator)
set(CMAKE_CXX_STANDARD 17)
set(CMAKE_CXX_STANDARD_REQUIRED ON)
find_package(Threads REQUIRED)
add_executable(RaftSimulator
    src/main.cpp
    src/raft_core/Node.cpp
)
target_include_directories(RaftSimulator PRIVATE
    ${CMAKE_CURRENT_SOURCE_DIR}/src
    ${CMAKE_CURRENT_SOURCE_DIR}/src/raft_core
    ${CMAKE_CURRENT_SOURCE_DIR}/src/telemetry
)
target_link_libraries(RaftSimulator PRIVATE Threads::Threads)
\end{lstlisting}

\begin{lstlisting}[style=bashstyle, caption={Direct g++ compilation (Windows)}]
g++ -I./src -I./asio-asio-1-28-0/asio/include -o raft_simulator.exe \
    src/main.cpp src/raft_core/Node.cpp -lpthread -lws2_32 -lwsock32
\end{lstlisting}

\newpage

% =====================================================================
\section{Step 2: Defining the Raft Data Model (Types.h)}
% =====================================================================

\subsection{Conceptual Overview}
The Raft protocol is a \textbf{replicated state machine}. Every node must
maintain an identical log of commands, and the consensus algorithm ensures
this log stays consistent across the cluster even during failures. Before
implementing any behavior, the data structures that represent this protocol
must be defined precisely. These types are the shared language that all
components of the system speak.

\subsection{The Building Process}
All protocol-level types are centralized in
\texttt{backend/src/raft\_core/Types.h}. This file defines no behavior---it
only defines \textit{what the data looks like}. This avoids circular
dependencies between \texttt{Node.h} and \texttt{main.cpp}.

\subsection{Implementation Details}
\begin{itemize}[leftmargin=*]
    \item \textbf{NodeState:} A scoped \texttt{enum class} with three
          values---\texttt{FOLLOWER}, \texttt{CANDIDATE}, \texttt{LEADER}.
          Scoped enums prevent name collisions and force explicit
          qualification (e.g., \texttt{NodeState::LEADER}).
    \item \textbf{LogEntry:} The fundamental unit of the replicated log.
          Stores the \texttt{term} in which it was created and the
          \texttt{command} string. The term is essential for Raft's safety
          proof: a leader can only commit entries from its current term.
    \item \textbf{RPC Structs:} \texttt{RequestVoteArgs/Reply} and
          \texttt{AppendEntriesArgs/Reply} define payload schemas for Raft's
          two core RPCs. In a real system these would be serialized to bytes;
          here they are passed directly as C++ structs between threads.
\end{itemize}

\subsection{Code \& Configuration}
\begin{lstlisting}[style=cppstyle, caption={backend/src/raft\_core/Types.h}]
#pragma once
#include <string>
#include <vector>
#include <optional>

namespace raft {

enum class NodeState { FOLLOWER, CANDIDATE, LEADER };

struct LogEntry {
    int term;
    std::string command;
};

struct RequestVoteArgs {
    int term; int candidateId;
    int lastLogIndex; int lastLogTerm;
};
struct RequestVoteReply { int term; bool voteGranted; };

struct AppendEntriesArgs {
    int term; int leaderId;
    int prevLogIndex; int prevLogTerm;
    std::vector<LogEntry> entries;
    int leaderCommit;
};
struct AppendEntriesReply { int term; bool success; };

} // namespace raft
\end{lstlisting}

\newpage

% =====================================================================
\section{Step 3: The Thread-Safe Telemetry Infrastructure}
% =====================================================================

\subsection{Conceptual Overview}
The C++ backend has two categories of threads running concurrently:
\textbf{Node threads} (run Raft logic, produce state-change events) and a
\textbf{Network I/O thread} (runs the Crow WebSocket server). These must
communicate without data races. The solution is the
\textbf{Producer-Consumer pattern}: nodes push events into a shared queue;
a dedicated consumer pops events and broadcasts them over WebSocket.

\subsection{The Building Process}
Two files define the infrastructure:
\begin{itemize}[leftmargin=*]
    \item \texttt{TelemetryTypes.h} --- the \texttt{RaftStateEvent} payload.
    \item \texttt{ThreadSafeQueue.h} --- a generic mutex-backed queue template.
\end{itemize}

\subsection{Implementation Details}
\textbf{ThreadSafeQueue} wraps \texttt{std::queue<T>} with a
\texttt{std::mutex} and \texttt{std::condition\_variable}.
\texttt{push()} uses \texttt{std::lock\_guard} (RAII: the mutex is
released on scope exit even if an exception is thrown) and calls
\texttt{m\_cond.notify\_one()} to wake exactly one sleeping consumer.

\texttt{wait\_and\_pop()} uses \texttt{std::unique\_lock} (can temporarily
release the lock). \texttt{m\_cond.wait(lock, predicate)} atomically
releases the lock and sleeps until both \texttt{notify\_one()} fires AND
the predicate \texttt{!m\_queue.empty()} is true, guarding against
\textbf{spurious wakeups}.

The copy constructor and assignment operator are explicitly
\texttt{= delete}d because \texttt{std::mutex} is non-copyable.

\subsection{Code \& Configuration}
\begin{lstlisting}[style=cppstyle, caption={backend/src/telemetry/ThreadSafeQueue.h}]
#pragma once
#include <queue>
#include <mutex>
#include <condition_variable>

namespace telemetry {

template <typename T>
class ThreadSafeQueue {
public:
    ThreadSafeQueue() = default;
    ~ThreadSafeQueue() = default;
    ThreadSafeQueue(const ThreadSafeQueue&) = delete;
    ThreadSafeQueue& operator=(const ThreadSafeQueue&) = delete;

    void push(T value) {
        std::lock_guard<std::mutex> lock(m_mutex);
        m_queue.push(std::move(value));
        m_cond.notify_one();
    }

    bool try_pop(T& value) {
        std::lock_guard<std::mutex> lock(m_mutex);
        if (m_queue.empty()) return false;
        value = std::move(m_queue.front());
        m_queue.pop();
        return true;
    }

    void wait_and_pop(T& value) {
        std::unique_lock<std::mutex> lock(m_mutex);
        m_cond.wait(lock, [this] { return !m_queue.empty(); });
        value = std::move(m_queue.front());
        m_queue.pop();
    }

private:
    std::queue<T>           m_queue;
    mutable std::mutex      m_mutex;
    std::condition_variable m_cond;
};

} // namespace telemetry
\end{lstlisting}

\begin{lstlisting}[style=cppstyle, caption={backend/src/telemetry/TelemetryTypes.h}]
#pragma once
#include "../raft_core/Types.h"
#include <string>
#include <vector>

namespace telemetry {

enum class EventType {
    STATE_TRANSITION, ELECTION_STARTED, VOTE_CAST
};

struct RaftStateEvent {
    int nodeId;
    EventType type;
    int currentTerm;
    raft::NodeState role;
    int commitIndex;
    std::string message;
    bool isCrashed = false;
    std::vector<raft::LogEntry> logEntries;
};

} // namespace telemetry
\end{lstlisting}

\newpage

% =====================================================================
\section{Step 4: The Node Class Interface (Node.h)}
% =====================================================================

\subsection{Conceptual Overview}
The \texttt{Node} class is the central abstraction of the entire system,
modeling a single server in the distributed cluster. Its public interface
defines what the outside world can do to a node, while private members
enforce strict encapsulation of internal Raft state.

\subsection{The Building Process}
\texttt{Node.h} declares the class with three state categories directly
mirroring the Raft paper:
\begin{itemize}[leftmargin=*]
    \item \textbf{Persistent state}: \texttt{m\_currentTerm},
          \texttt{m\_votedFor}, \texttt{m\_log} (must survive crashes).
    \item \textbf{Volatile state}: \texttt{m\_commitIndex},
          \texttt{m\_lastApplied}, \texttt{m\_state} (reconstructable).
    \item \textbf{Volatile leader state}: \texttt{m\_nextIndex},
          \texttt{m\_matchIndex} (initialized fresh on each election win).
\end{itemize}

\subsection{Implementation Details}
\begin{itemize}[leftmargin=*]
    \item \textbf{Atomic booleans}: \texttt{m\_running} and
          \texttt{m\_isCrashed} are \texttt{std::atomic<bool>}, supporting
          lock-free, indivisible read-modify-write operations across threads.
    \item \textbf{std::optional}: \texttt{m\_votedFor} is
          \texttt{std::optional<int>}---a node either has voted for someone
          or has not voted (\texttt{std::nullopt}). More expressive and
          safer than a \texttt{-1} sentinel.
    \item \textbf{Dependency Injection}: RPC dispatchers are
          \texttt{std::function} objects injected from \texttt{main.cpp}.
          The node does not know \textit{how} RPCs are delivered, cleanly
          separating the algorithm from the networking layer.
\end{itemize}

\subsection{Code \& Configuration}
\begin{lstlisting}[style=cppstyle, caption={backend/src/raft\_core/Node.h}]
#pragma once
#include "Types.h"
#include "../telemetry/TelemetryTypes.h"
#include <mutex>
#include <condition_variable>
#include <thread>
#include <atomic>
#include <vector>
#include <chrono>
#include <functional>
#include <map>

namespace raft {

class Node {
public:
    Node(int id, const std::vector<int>& peerIds);
    ~Node();
    void start(); void stop();
    RequestVoteReply requestVote(const RequestVoteArgs& args);
    AppendEntriesReply appendEntries(const AppendEntriesArgs& args);
    void setRpcHandler(
        std::function<RequestVoteReply(int, const RequestVoteArgs&)>,
        std::function<AppendEntriesReply(int, const AppendEntriesArgs&)>
    );
    void setTelemetryCallback(
        std::function<void(const telemetry::RaftStateEvent&)> cb);
    void forceSyncTelemetry();
    void crash(); void recover();
    void clientRequest(const std::string& command);

private:
    void emitTelemetry(telemetry::EventType type,
                       const std::string& message = "");
    void run(); void resetElectionTimer();
    void startElection(); void becomeFollower(int term);
    void becomeLeader();

    int m_id; std::vector<int> m_peerIds;
    std::atomic<bool> m_running{false};
    std::atomic<bool> m_isCrashed{false};
    std::thread m_workerThread;
    std::mutex m_mutex; std::condition_variable m_cv;

    // Persistent Raft state
    int m_currentTerm = 0;
    std::optional<int> m_votedFor;
    std::vector<LogEntry> m_log;

    // Volatile Raft state
    int m_commitIndex = 0; int m_lastApplied = 0;
    NodeState m_state = NodeState::FOLLOWER;

    // Volatile leader state
    std::map<int,int> m_nextIndex; std::map<int,int> m_matchIndex;

    std::chrono::time_point<std::chrono::steady_clock> m_lastHeartbeat;
    std::chrono::milliseconds m_electionTimeout;

    std::function<RequestVoteReply(int, const RequestVoteArgs&)>
        sendRequestVote;
    std::function<AppendEntriesReply(int, const AppendEntriesArgs&)>
        sendAppendEntries;
    std::function<void(const telemetry::RaftStateEvent&)>
        m_telemetryCallback;
};

} // namespace raft
\end{lstlisting}

\newpage

% =====================================================================
\section{Step 5: The Node Lifecycle---Run Loop and Leader Election}
% =====================================================================

\subsection{Conceptual Overview}
A Raft node starts as a \textbf{Follower}. If it hears nothing within its
election timeout, it promotes itself to \textbf{Candidate} and launches an
election. If it wins (receives votes from a strict majority), it becomes
\textbf{Leader}. The quorum requirement makes this safe:
\begin{equation}
    Q = \left\lfloor \frac{N}{2} \right\rfloor + 1
\end{equation}
Any two majorities overlap by at least one node, making it
\textit{mathematically impossible} for two leaders to be elected in the
same term.

\subsection{The Building Process}
Three methods implement this lifecycle in \texttt{Node.cpp}:
\begin{enumerate}[leftmargin=*]
    \item \texttt{run()} -- the main loop; checks the election timer.
    \item \texttt{resetElectionTimer()} -- randomizes the timeout.
    \item \texttt{startElection()} -- transitions to Candidate and
          dispatches concurrent \texttt{RequestVote} RPCs.
\end{enumerate}

\subsection{Implementation Details}
\begin{itemize}[leftmargin=*]
    \item \textbf{Randomized timeouts:} \texttt{thread\_local std::mt19937}
          gives each thread its own independent Mersenne Twister PRNG,
          avoiding contention on a global object. Timeouts range 150--300ms.
    \item \textbf{Fire-and-forget RPCs:} A new \texttt{std::thread} is
          launched per peer and immediately \texttt{.detach()}ed. The RPC
          call blocks \textit{inside} the thread while the node thread stays
          responsive.
    \item \textbf{Stale state protection:} After a detached thread receives
          a reply, it re-acquires the mutex and verifies the node is still a
          Candidate for the same term before processing the vote.
    \item \textbf{Shared vote counter:} \texttt{std::make\_shared<int>(1)}
          captured by value lets multiple concurrent threads safely share
          and increment the counter once they hold the lock.
\end{itemize}

\subsection{Code \& Configuration}
\begin{lstlisting}[style=cppstyle, caption={Node.cpp --- run() and resetElectionTimer()}]
void Node::run() {
    while (m_running) {
        std::unique_lock<std::mutex> lock(m_mutex);
        auto now = std::chrono::steady_clock::now();
        if (!m_isCrashed && m_state != NodeState::LEADER
            && now - m_lastHeartbeat >= m_electionTimeout)
        {
            startElection();
        }
        m_cv.wait_for(lock, std::chrono::milliseconds(50),
            [this]() { return !m_running; });
    }
}

void Node::resetElectionTimer() {
    static thread_local std::mt19937
        generator(std::random_device{}());
    std::uniform_int_distribution<int> dist(150, 300);
    m_electionTimeout = std::chrono::milliseconds(dist(generator));
    m_lastHeartbeat = std::chrono::steady_clock::now();
}
\end{lstlisting}

\begin{lstlisting}[style=cppstyle, caption={Node.cpp --- startElection()}]
void Node::startElection() {
    m_state = NodeState::CANDIDATE;
    m_currentTerm++; m_votedFor = m_id;
    resetElectionTimer();

    RequestVoteArgs args;
    args.term        = m_currentTerm; args.candidateId = m_id;
    args.lastLogIndex = m_log.empty() ? 0 : (int)m_log.size();
    args.lastLogTerm  = m_log.empty() ? 0 : m_log.back().term;

    auto votesGranted = std::make_shared<int>(1);
    int majority = (m_peerIds.size() + 1) / 2;

    for (int peerId : m_peerIds) {
        std::thread([this, peerId, args, votesGranted, majority]() {
            RequestVoteReply reply = sendRequestVote(peerId, args);

            std::lock_guard<std::mutex> lock(m_mutex);
            if (m_state != NodeState::CANDIDATE
                || m_currentTerm != args.term) return; // Stale

            if (reply.term > m_currentTerm) {
                becomeFollower(reply.term); return;
            }
            if (reply.voteGranted) {
                (*votesGranted)++;
                if (*votesGranted > majority) becomeLeader();
            }
        }).detach();
    }
}
\end{lstlisting}

\newpage

% =====================================================================
\section{Step 6: The RequestVote and AppendEntries RPCs}
% =====================================================================

\subsection{Conceptual Overview}
The two RPCs are the only means of inter-node communication in Raft.
\textbf{RequestVote} is called during elections. \textbf{AppendEntries}
serves a dual purpose: replicating log entries and acting as a periodic
heartbeat that suppresses elections. These two functions collectively encode
the entire Raft safety model.

\subsection{Implementation Details}
\textbf{RequestVote receiver:}
\begin{enumerate}[leftmargin=*]
    \item Refuse if crashed (simulates a severed connection).
    \item Refuse if the candidate's term is stale.
    \item If the candidate's term is newer, step down to Follower (Term
          Ascendancy Rule).
    \item Grant the vote only if we haven't voted yet AND the candidate's
          log is at least as up-to-date as ours (Leader Completeness
          Property).
\end{enumerate}

\textbf{AppendEntries receiver:}
\begin{enumerate}[leftmargin=*]
    \item Refuse if crashed or if sender's term is stale.
    \item Step down to Follower and reset the election timer (heartbeat).
    \item Perform the \textbf{Log Consistency Check}: verify log contains
          an entry at \texttt{prevLogIndex} with term \texttt{prevLogTerm}.
          If not, return failure so the leader retries with an older index.
    \item Append new entries, first deleting any conflicting ones.
    \item Advance \texttt{m\_commitIndex} if \texttt{leaderCommit} is greater.
\end{enumerate}

\subsection{Code \& Configuration}
\begin{lstlisting}[style=cppstyle, caption={Node.cpp --- requestVote()}]
RequestVoteReply Node::requestVote(const RequestVoteArgs& args) {
    std::lock_guard<std::mutex> lock(m_mutex);
    RequestVoteReply reply;
    reply.term = m_currentTerm; reply.voteGranted = false;

    if (m_isCrashed)               return reply;
    if (args.term < m_currentTerm) return reply;
    if (args.term > m_currentTerm) becomeFollower(args.term);

    if (!m_votedFor.has_value()
        || m_votedFor.value() == args.candidateId)
    {
        int lastLogTerm  = m_log.empty() ? 0 : m_log.back().term;
        int lastLogIndex = m_log.empty() ? 0 : (int)m_log.size();
        bool logIsUpToDate =
            (args.lastLogTerm > lastLogTerm) ||
            (args.lastLogTerm == lastLogTerm
             && args.lastLogIndex >= lastLogIndex);
        if (logIsUpToDate) {
            m_votedFor = args.candidateId;
            reply.voteGranted = true;
            resetElectionTimer();
        }
    }
    reply.term = m_currentTerm;
    return reply;
}
\end{lstlisting}

\begin{lstlisting}[style=cppstyle, caption={Node.cpp --- appendEntries()}]
AppendEntriesReply Node::appendEntries(const AppendEntriesArgs& args){
    std::lock_guard<std::mutex> lock(m_mutex);
    AppendEntriesReply reply;
    reply.term = m_currentTerm; reply.success = false;

    if (m_isCrashed)               return reply;
    if (args.term < m_currentTerm) return reply;
    if (args.term > m_currentTerm || m_state != NodeState::FOLLOWER)
        becomeFollower(args.term);

    resetElectionTimer(); // Leader is alive

    if (args.prevLogIndex > (int)m_log.size()) return reply;
    if (args.prevLogIndex > 0
        && m_log[args.prevLogIndex-1].term != args.prevLogTerm)
        return reply;

    reply.success = true; reply.term = m_currentTerm;

    int index = args.prevLogIndex;
    for (const auto& entry : args.entries) {
        index++;
        if ((int)m_log.size() >= index) {
            if (m_log[index-1].term != entry.term) {
                m_log.erase(m_log.begin()+index-1, m_log.end());
                m_log.push_back(entry);
            }
        } else { m_log.push_back(entry); }
    }

    if (args.leaderCommit > m_commitIndex)
        m_commitIndex = std::min(args.leaderCommit,(int)m_log.size());

    return reply;
}
\end{lstlisting}

\newpage

% =====================================================================
\section{Step 7: Leader Heartbeats, Log Replication, and Commitment}
% =====================================================================

\subsection{Conceptual Overview}
After winning an election, the leader must broadcast \texttt{AppendEntries}
heartbeats every 50ms to suppress follower timeouts. When a client mutation
arrives, the leader appends it to its own log and replicates it. Once a
majority quorum have written the entry, the leader advances its
\texttt{commitIndex} and the command is durably committed.

\subsection{Implementation Details}
\begin{itemize}[leftmargin=*]
    \item \textbf{Term-pinned heartbeat thread:} The background thread
          captures \texttt{termAtElection} and exits if the term changes,
          preventing stale threads from interfering with a new leadership.
    \item \textbf{nextIndex / matchIndex:} \texttt{m\_nextIndex[peer]}
          tracks the next log index to send; \texttt{m\_matchIndex[peer]}
          tracks the highest index confirmed replicated. On failure, the
          leader decrements \texttt{nextIndex} and retries.
    \item \textbf{Quorum commit:} The leader scans backwards for the highest
          index $n$ where \texttt{log[n-1].term == currentTerm} and a
          majority of nodes have \texttt{matchIndex $\geq$ n}.
\end{itemize}

\subsection{Code \& Configuration}
\begin{lstlisting}[style=cppstyle, caption={Node.cpp --- becomeLeader() and clientRequest()}]
void Node::becomeLeader() {
    m_state = NodeState::LEADER;
    for (int p : m_peerIds) {
        m_nextIndex[p]  = m_log.size() + 1;
        m_matchIndex[p] = 0;
    }

    std::thread([this, termAtElection = m_currentTerm]() {
        while (m_running) {
            {
                std::lock_guard<std::mutex> lock(m_mutex);
                if (m_state!=NodeState::LEADER
                    || m_currentTerm!=termAtElection) break;
                if (m_isCrashed) continue;

                for (int peerId : m_peerIds) {
                    AppendEntriesArgs args;
                    args.term = m_currentTerm; args.leaderId = m_id;
                    args.leaderCommit = m_commitIndex;
                    int nextIdx = m_nextIndex[peerId];
                    args.prevLogIndex = nextIdx - 1;
                    args.prevLogTerm =
                        (args.prevLogIndex > 0
                         && args.prevLogIndex <= (int)m_log.size())
                        ? m_log[args.prevLogIndex-1].term : 0;
                    if (nextIdx>0 && nextIdx<=(int)m_log.size())
                        args.entries = std::vector<LogEntry>(
                            m_log.begin()+nextIdx-1, m_log.end());
                    else args.entries = {};

                    std::thread([this, peerId, args]() {
                        AppendEntriesReply reply =
                            sendAppendEntries(peerId, args);
                        std::lock_guard<std::mutex> lock(m_mutex);
                        if (m_state!=NodeState::LEADER
                            || m_currentTerm!=args.term) return;
                        if (reply.term > m_currentTerm) {
                            becomeFollower(reply.term); return;
                        }
                        if (reply.success) {
                            int mi = args.prevLogIndex
                                    +(int)args.entries.size();
                            if (mi > m_matchIndex[peerId]) {
                                m_matchIndex[peerId] = mi;
                                m_nextIndex[peerId]  = mi + 1;
                                for (int n=(int)m_log.size();
                                     n>m_commitIndex; --n) {
                                    if(m_log[n-1].term==m_currentTerm){
                                        int cnt=1;
                                        for(int p:m_peerIds)
                                            if(m_matchIndex[p]>=n)cnt++;
                                        if(cnt>(int)(m_peerIds.size()+1)/2){
                                            m_commitIndex=n; break;
                                        }
                                    }
                                }
                            }
                        } else {
                            m_nextIndex[peerId] =
                                std::max(1,m_nextIndex[peerId]-1);
                        }
                    }).detach();
                }
            }
            std::this_thread::sleep_for(std::chrono::milliseconds(50));
        }
    }).detach();
}

void Node::clientRequest(const std::string& command) {
    std::lock_guard<std::mutex> lock(m_mutex);
    if (m_state != NodeState::LEADER || m_isCrashed) return;
    m_log.push_back({m_currentTerm, command});
    emitTelemetry(telemetry::EventType::STATE_TRANSITION,
                  "Appended client request");
}
\end{lstlisting}

\newpage

% =====================================================================
\section{Step 8: The Orchestrator---main.cpp}
% =====================================================================

\subsection{Conceptual Overview}
\texttt{main.cpp} is the glue layer: it instantiates all nodes, wires them
together via Dependency Injection, launches the telemetry consumer thread,
and runs the Crow WebSocket server. It contains no Raft algorithm logic.

\subsection{Implementation Details}
\begin{itemize}[leftmargin=*]
    \item \textbf{RPC wiring:} Lambdas capture \texttt{\&nodes} by reference
          and call \texttt{simulateNetworkLatency()} before dispatching to
          \texttt{nodes[targetId-1]->requestVote(args)}---a direct function
          call across a thread boundary simulating a network round-trip.
    \item \textbf{onopen:} Inserts the new connection pointer and calls
          \texttt{forceSyncTelemetry()} so the new client immediately sees
          full cluster state.
    \item \textbf{onmessage routing:} \texttt{CLIENT\_REQUEST} is broadcast
          to all nodes (only the leader processes it internally).
          \texttt{CRASH\_NODE} / \texttt{RECOVER\_NODE} target a specific
          node by integer ID.
\end{itemize}

\subsection{Code \& Configuration}
\begin{lstlisting}[style=cppstyle, caption={backend/src/main.cpp --- Orchestrator (abridged)}]
int main() {
    telemetry::ThreadSafeQueue<telemetry::RaftStateEvent> queue;

    const int N = 5;
    std::vector<std::shared_ptr<Node>> nodes;
    for (int i = 1; i <= N; ++i) {
        std::vector<int> peers;
        for (int j = 1; j <= N; ++j) if (i!=j) peers.push_back(j);
        nodes.push_back(std::make_shared<Node>(i, peers));
    }

    for (int i = 0; i < N; ++i) {
        nodes[i]->setRpcHandler(
            [&nodes](int id, const RequestVoteArgs& a) {
                simulateNetworkLatency();
                return nodes[id-1]->requestVote(a); },
            [&nodes](int id, const AppendEntriesArgs& a) {
                simulateNetworkLatency();
                return nodes[id-1]->appendEntries(a); }
        );
        nodes[i]->setTelemetryCallback(
            [&queue](const telemetry::RaftStateEvent& ev) {
                queue.push(ev); }
        );
    }

    std::mutex mtx;
    std::unordered_set<crow::websocket::connection*> clients;

    std::thread([&]() {
        while (true) {
            telemetry::RaftStateEvent ev;
            queue.wait_and_pop(ev);
            // ... serialize ev to JSON string payload ...
            std::lock_guard<std::mutex> lock(mtx);
            for (auto* c : clients) c->send_text(payload);
        }
    }).detach();

    for (auto& n : nodes) n->start();

    crow::SimpleApp app;
    CROW_WEBSOCKET_ROUTE(app, "/raft")
        .onopen([&](crow::websocket::connection& c) {
            std::lock_guard<std::mutex> lock(mtx);
            clients.insert(&c);
            for (auto& n : nodes) n->forceSyncTelemetry();
        })
        .onclose([&](crow::websocket::connection& c, const std::string&){
            std::lock_guard<std::mutex> lock(mtx);
            clients.erase(&c);
        })
        .onmessage([&nodes](crow::websocket::connection&,
                            const std::string& data, bool) {
            auto p = crow::json::load(data);
            if (!p) return;
            if (p.has("command")) {
                std::string cmd = p["command"].s();
                if (cmd=="CLIENT_REQUEST" && p.has("data"))
                    for (auto& n : nodes)
                        n->clientRequest(p["data"].s());
                else if (p.has("nodeId")) {
                    int id = p["nodeId"].i();
                    if (id>=1 && id<=(int)nodes.size()) {
                        if (cmd=="CRASH_NODE")   nodes[id-1]->crash();
                        if (cmd=="RECOVER_NODE") nodes[id-1]->recover();
                    }
                }
            }
        });

    app.port(8080).multithreaded().run();
    return 0;
}
\end{lstlisting}

\newpage

% =====================================================================
\section{Step 9: Redux State Management and the WebSocket Hook}
% =====================================================================

\subsection{Conceptual Overview}
On the frontend, the challenge is managing a stream of high-frequency,
partial state updates and reflecting them in a reactive UI. Redux Toolkit's
\texttt{createSlice} provides an immutable, globally accessible state
container. A custom Hook (\texttt{useRaftWebSocket}) manages the WebSocket
lifecycle and feeds incoming messages into the Redux store.

\subsection{Implementation Details}
\begin{itemize}[leftmargin=*]
    \item \textbf{Redux Toolkit / Immer:} The reducer appears to mutate
          state directly. Redux Toolkit wraps reducers with Immer, which
          intercepts assignments on a Proxy draft and produces a new
          immutable state tree transparently.
    \item \textbf{useRef for the socket:} Updating a ref does \textit{not}
          trigger a re-render. This is critical because the WebSocket
          transitions through internal states (CONNECTING, OPEN, CLOSING,
          CLOSED) frequently.
    \item \textbf{Cleanup function:} The function returned from
          \texttt{useEffect} closes the socket on unmount, preventing zombie
          connections from calling \texttt{dispatch} on a destroyed component.
    \item \textbf{Polymorphic sendCommand:} Inspects \texttt{typeof payload}
          at runtime to route either a \texttt{nodeId} number or a
          \texttt{data} string, serving all UI panels from one function.
\end{itemize}

\subsection{Code \& Configuration}
\begin{lstlisting}[style=tsstyle, caption={frontend/src/store/raftSlice.ts}]
import { createSlice, PayloadAction } from '@reduxjs/toolkit';

export interface LogEntry { term: number; command: string; }
export interface RaftNodeState {
  nodeId: number; term: number;
  role: number; commitIndex: number;
  isCrashed: boolean; log: LogEntry[];
}
interface RaftPayload {
  nodeId: number; term: number; role: number;
  commitIndex: number; isCrashed?: boolean; log?: LogEntry[];
}

export const raftSlice = createSlice({
  name: 'raft',
  initialState: {
    nodes: {
      1:{nodeId:1,term:0,role:0,commitIndex:0,isCrashed:false,log:[]},
      2:{nodeId:2,term:0,role:0,commitIndex:0,isCrashed:false,log:[]},
      3:{nodeId:3,term:0,role:0,commitIndex:0,isCrashed:false,log:[]},
      4:{nodeId:4,term:0,role:0,commitIndex:0,isCrashed:false,log:[]},
      5:{nodeId:5,term:0,role:0,commitIndex:0,isCrashed:false,log:[]},
    } as Record<number, RaftNodeState>
  },
  reducers: {
    updateNodeState:(state, action: PayloadAction<RaftPayload>) => {
      const { nodeId, term, role, commitIndex, isCrashed, log }
        = action.payload;
      if (state.nodes[nodeId]) {
        state.nodes[nodeId].term        = term;
        state.nodes[nodeId].role        = role;
        state.nodes[nodeId].commitIndex = commitIndex;
        if (isCrashed!==undefined)
          state.nodes[nodeId].isCrashed = isCrashed;
        if (log!==undefined)
          state.nodes[nodeId].log = log;
      }
    }
  }
});
\end{lstlisting}

\begin{lstlisting}[style=tsstyle, caption={frontend/src/hooks/useRaftWebSocket.ts}]
import { useEffect, useRef } from 'react';
import { useDispatch } from 'react-redux';
import { updateNodeState } from '../store/raftSlice';

export const useRaftWebSocket = (url: string) => {
  const dispatch = useDispatch();
  const ws = useRef<WebSocket | null>(null);

  useEffect(() => {
    ws.current = new WebSocket(url);
    ws.current.onmessage = (event) => {
      try { dispatch(updateNodeState(JSON.parse(event.data))); }
      catch (e) { console.error('Parse error', e); }
    };
    return () => {
      if (ws.current?.readyState === WebSocket.OPEN)
        ws.current.close();
    };
  }, [url, dispatch]);

  const sendCommand = (command: string, payload?: any) => {
    if (ws.current?.readyState !== WebSocket.OPEN) return;
    if (typeof payload === 'number')
      ws.current.send(JSON.stringify({ command, nodeId: payload }));
    else if (typeof payload === 'string')
      ws.current.send(JSON.stringify({ command, data: payload }));
    else
      ws.current.send(JSON.stringify({ command }));
  };

  return { sendCommand };
};
\end{lstlisting}

\newpage

% =====================================================================
\section{Step 10: React UI---Visualization and Chaos Engineering}
% =====================================================================

\subsection{Conceptual Overview}
The final layer transforms Redux state into a living, interactive
visualization. \texttt{ClusterView} orchestrates the layout;
\texttt{NodeCard} visualizes each node; \texttt{ChaosControlPanel} and
\texttt{ClientPanel} enable interactive fault injection and log mutations.
Every state change in the C++ backend appears in the UI within one 50ms
heartbeat cycle.

\subsection{The Building Process}
Component hierarchy:
\begin{center}
\texttt{App} $\rightarrow$ \texttt{ClusterView} $\rightarrow$
\{\texttt{NodeCard}$\times$5, \texttt{ClientPanel}, \texttt{ChaosControlPanel}\}
\end{center}
Only \texttt{ClusterView} calls \texttt{useRaftWebSocket}, obtaining
\texttt{sendCommand} and passing it as props to child panels.

\subsection{Implementation Details}
\begin{itemize}[leftmargin=*]
    \item \textbf{Role-based styling:} \texttt{NodeCard} maps
          \texttt{node.role} (0/1/2) to distinct Tailwind CSS class sets.
          Crashed nodes receive a desaturated red style.
    \item \textbf{Log display:} \texttt{node.log.slice(-2)} provides a
          rolling two-entry window without overflowing the card height.
    \item \textbf{ChaosControlPanel:} Renders Crash/Recover buttons per node.
          Buttons are disabled based on \texttt{node.isCrashed} to prevent
          double-crashing.
    \item \textbf{ClientPanel:} A controlled form that dispatches
          \texttt{CLIENT\_REQUEST} and resets its input on submit.
\end{itemize}

\subsection{Code \& Configuration}
\begin{lstlisting}[style=tsstyle, caption={frontend/src/components/NodeCard.tsx}]
import React from 'react';
import type { RaftNodeState } from '../store/raftSlice';

export const NodeCard: React.FC<{node: RaftNodeState}> = ({ node }) => {
  const roleStyles: Record<number,string> = {
    0: 'bg-slate-800 border-slate-600 text-slate-300',
    1: 'bg-amber-600 border-amber-400 text-white',
    2: 'bg-blue-600 border-blue-400 text-white shadow-[0_0_20px_rgba(37,99,235,0.7)]'
  };
  const roleNames: Record<number,string> = {
    0:'FOLLOWER', 1:'CANDIDATE', 2:'LEADER'
  };
  const currentStyle = node.isCrashed
    ? 'bg-red-950 border-red-800 text-red-500 opacity-60'
    : roleStyles[node.role];
  const roleName = node.isCrashed
    ? 'CRASHED' : roleNames[node.role];

  return (
    <div className={`flex flex-col p-6 m-3 rounded-2xl border-2
                     transition-all duration-300 ${currentStyle}`}>
      <h2 className="text-3xl font-bold mb-2">Node {node.nodeId}</h2>
      <div className="tracking-widest mb-5">{roleName}</div>
      <div className="font-mono space-y-2">
        <p>Term: {node.term}</p>
        <p>Commit Idx: {node.commitIndex}</p>
        {node.log.length > 0 && (
          <ul>
            {node.log.slice(-2).map((e, i) => (
              <li key={i}>
                <span>T{e.term} </span>
                <span className="text-emerald-300">{e.command}</span>
              </li>
            ))}
          </ul>
        )}
      </div>
    </div>
  );
};
\end{lstlisting}

\begin{lstlisting}[style=tsstyle, caption={frontend/src/components/ClientPanel.tsx}]
import React, { useState } from 'react';

interface ClientPanelProps {
  onSendCommand: (command: string, data?: string) => void;
}

export const ClientPanel: React.FC<ClientPanelProps> = ({onSendCommand}) => {
  const [inputValue, setInputValue] = useState('');

  const handleSubmit = (e: React.FormEvent) => {
    e.preventDefault();
    if (inputValue.trim()) {
      onSendCommand('CLIENT_REQUEST', inputValue.trim());
      setInputValue('');
    }
  };

  return (
    <form onSubmit={handleSubmit} className="flex gap-4">
      <input
        type="text"
        value={inputValue}
        onChange={(e) => setInputValue(e.target.value)}
        placeholder="e.g. SET X=10"
        className="flex-1 px-4 py-2 bg-slate-900 border rounded text-white"
      />
      <button type="submit" disabled={!inputValue.trim()}
              className="px-6 py-2 bg-blue-600 disabled:bg-slate-600
                         text-white font-bold rounded">
        Send to Leader
      </button>
    </form>
  );
};
\end{lstlisting}

\section*{Codebase Modifications: Cloud Deployment}
To deploy the Raft Consensus Simulator to Render (backend) and Vercel (frontend), the following architectural changes were made:
\begin{itemize}
    \item \textbf{Backend Dockerization:} Created a \texttt{Dockerfile} for the C++ backend using \texttt{ubuntu:22.04} to compile and serve the binary, enabling deployment as a Web Service on Render since Render does not provide a native C++ runtime.
    \item \textbf{Port Binding Update:} Modified \texttt{backend/src/main.cpp} to dynamically read the \texttt{PORT} environment variable (defaulting to 8080) and bind the Crow WebSocket server to \texttt{0.0.0.0} instead of \texttt{localhost}. This is required for Render to properly route external traffic to the container.
    \item \textbf{Frontend Environment Variables:} Updated \texttt{frontend/src/components/ClusterView.tsx} to read the WebSocket URL from Vite environment variables (\texttt{import.meta.env.VITE\_WS\_URL}) rather than hardcoding \texttt{ws://localhost:8080/raft}. This allows Vercel to dynamically point the frontend to the deployed Render backend in production.
\end{itemize}

\end{document}
